% !PRINTED 639
\weberpage{639}
9. \gesperrt{A system of integral functions derived from a minimal
 basis $\omega_1, \omega_2, \ldots, \omega_n$}

\weq{7}{\omega_\nu' = x_{\nu,1} \omega_1 + x_{\nu,2} \omega_2 + \cdots + x_{\nu,n} \omega_n,}
\[
\nu = 1, 2, \ldots, n
\]

\gesperrt{is then and only then a minimal basis, when the
determinant}

\weq{8}{X = \Sigma \pm x_{1,1} x_{2,2} \ldots x_{n,n}}

\gesperrt{is a constant different from zero.}

For if this determinant, which is a rational function
of $z$, has a linear factor $z - c$, then the constants
$c_1, c_2, \ldots, c_n$ can be so determined that the $n$ integral rational functions


\[
c_1 x_{1,\nu} + c_2 x_{2,\nu} + \cdots + c_n x_{n,\nu}
\]

become divisible by $z - c$ without all the $c_i$ vanishing.
But then

\[
\frac{c_1 \omega_1' + c_2 \omega_2' + \cdots + c_n \omega_n'}{z - c}
\]

is an integral function, and $\omega_1', \omega_2', \ldots, \omega_n'$ is not a minimal basis.

For the discriminant one obtains from (7), according to § 173, (6)

\weq{9}{\varDelta(\omega_1', \omega_2', \ldots, \omega_n') = X^2 \varDelta(\omega_1, \omega_2, \ldots, \omega_n).}

From this it follows:

10. \gesperrt{The discriminant of a minimal basis is, apart from
a constant factor, independent of the particular
 choice of the basis.}

Among all discriminants of bases formed from integral functions,
 the discriminant of a minimal basis has the \gesperrt{lowest
 degree} (hence the name minimal basis). It is an integral
rational function of $z$ uniquely determined by the field
itself, apart from a constant factor. It is therefore also called
\gesperrt{the discriminant of the field $\Omega$} and is denoted
by $\varDelta(\Omega)$.
